\section{Methodology}

Our approach proceeds in two phases: (1) compression response profiling to generate a task-configuration matrix, and (2) statistical analysis to determine whether tasks cluster and whether attention features predict cluster membership.

\subsection{Compression Response Profiling}

We define 6 compression configurations spanning eviction and quantization dimensions:

\begin{table}[h]
\centering
\begin{tabular}{@{}llll@{}}
\toprule
Config & Method & Retention & Quantization \\
\midrule
C1 & Full & 100\% & FP16 \\
C2 & H2O & 80\% & FP16 \\
C3 & H2O & 40\% & FP16 \\
C4 & Full & 100\% & INT8 \\
C5 & Full & 100\% & INT4 \\
C6 & H2O & 60\% & INT8 \\
\bottomrule
\end{tabular}
\end{table}

For each of 21 LongBench tasks $t$ and each configuration $c$, we compute accuracy $A(t, c)$ using task-appropriate metrics. The response matrix $\mathbf{R} \in \mathbb{R}^{21 \times 6}$ contains accuracy retention ratios:

\begin{equation}
R_{tc} = \frac{A(t, c)}{A(t, c_{\text{full}})}
\end{equation}

\subsection{Cluster Discovery via Gap Statistic}

We use the gap statistic \citep{tibshirani2001gap} to determine the optimal number of clusters $k^*$ directly from data, avoiding bias from assuming task-compression structure exists.

For $k = 1, \ldots, K_{\max}$:
\begin{enumerate}
    \item Cluster $\mathbf{R}$ into $k$ groups using $k$-means, computing within-cluster dispersion $W_k$
    \item Generate $B$ reference datasets from uniform distribution on the same bounding box
    \item Compute $\text{Gap}(k) = E^*[\log W_k] - \log W_k$
\end{enumerate}

The optimal $k^*$ is the smallest $k$ satisfying: $\text{Gap}(k) \geq \text{Gap}(k+1) - s_{k+1}$

\subsection{Attention Entropy Analysis}

For each task, we compute attention entropy from the first 100 tokens across all 32 layers and 32 heads:

\begin{equation}
H(\alpha) = -\sum_{i=1}^{n} \alpha_i \log(\alpha_i + \epsilon)
\end{equation}

We perform one-way ANOVA across 6 LongBench task categories with entropy as the dependent variable.
